% Нарративная часть заключения (подключена в conclusion.tex через \input).
% Идёт перед "Основные результаты работы заключаются в следующем." (%% Согласно ГОСТ Р 7.0.11-2011:
%% 5.3.3 В заключении диссертации излагают итоги выполненного исследования, рекомендации, перспективы дальнейшей разработки темы.
%% 9.2.3 В заключении автореферата диссертации излагают итоги данного исследования, рекомендации и перспективы дальнейшей разработки темы.
\begin{enumerate}
  \item Нелинейное снижение размерности (лапласовы собственные карты, автокодировщики) восстанавливает геометрию среды из популяционной активности гиппокампа при свободном поведении, тогда как линейные методы (PCA) не справляются; качество восстановления растёт с опытом животного и числом зарегистрированных нейронов.

  \item Разработан алгоритм оценки эффективной размерности популяционной активности с коррекцией спектра корреляционных матриц для конечных записей; эффективная размерность снижается во время движения животного. Внутренняя размерность той же активности почти постоянна во времени, значимо ниже эффективной и демонстрирует степенное масштабирование с числом нейронов, что подтверждает нелинейную искривлённость многообразия.

  \item Предложен мультиплексный граф рекуррентности популяции~--- способ агрегации индивидуальных рекуррентных графов нейронов в единое представление, впервые применённый к сырым сигналам кальциевого имиджинга \textit{in~vivo}; из одного и того же графа восстанавливаются как геометрия среды (медианная ошибка декодирования положения около 9~см при стороне арены 44~см), так и динамические режимы коллективной активности, связанные со скоростью движения.

  \item Те же методы анализа применены к искусственным рекуррентным нейронным сетям: установлено, что при равной поведенческой эффективности обучение с подкреплением и обучение с учителем формируют качественно различные вычислительные архитектуры.

  \item Разработан информационно-теоретический фреймворк INTENSE для выявления и количественной оценки нейронной селективности в данных кальциевого имиджинга; показано, что около трети наблюдаемой мультиселективности обусловлено корреляциями поведенческих переменных, а замена взаимной информации на линейную корреляцию приводит к потере более двух третей значимых ассоциаций.

  \item Разработан фреймворк MANGO для безградиентной максимизации активации нейронов импульсных нейронных сетей (оптимизация с разложением тензорного поезда, в 3--15 раз быстрее эталонных методов); обнаружена двухфазная динамика формирования селективности, а последующий анализ инструментами DRIADA и INTENSE выявил двухфазную траекторию в информационной плоскости глубоких слоёв и усиление связи индивидуальной селективности с глубиной сети.

  \item Все разработанные методы объединены в открытый программный пакет DRIADA, интегрирующий анализ нейронной активности на индивидуальном, популяционном и сетевом уровнях; автоматический контроль качества выделенных компонент кальциевого имиджинга (Autoinspect) обеспечил воспроизводимую предобработку записей для всех нейрофизиологических разделов работы.
\end{enumerate}
).

В настоящей работе разработаны и применены методы анализа скрытых переменных коллективной динамики нейронной активности на нескольких взаимодополняющих уровнях: извлечение коллективных координат методами снижения размерности, количественная характеристика нейронного многообразия и анализ индивидуальной нейронной селективности.

На уровне коллективных переменных показано, что нелинейное снижение размерности (лапласовы собственные карты) позволяет восстановить геометрию среды из популяционной активности нейронов гиппокампа мышей при свободном исследовании, тогда как линейные методы (PCA) не справляются с данной задачей. Качество восстановления улучшается по мере знакомства животного со средой и растёт с числом зарегистрированных нейронов, что указывает на формирование когерентного популяционного кода. Те же аналитические подходы были применены к искусственным рекуррентным нейронным сетям, обученным задаче контекстно-зависимого принятия решений. Установлено, что при одинаковой поведенческой эффективности обучение с подкреплением и обучение с учителем формируют качественно различные вычислительные архитектуры: первое порождает гибридную динамику с квазипериодическими аттракторами и сбалансированными функциональными популяциями, тогда как второе предпочитает более простые решения на основе неподвижных точек.

Для количественной характеристики нейронного многообразия разработаны методы оценки его размерности. Предложен алгоритм расчёта эффективной размерности с коррекцией спектров корреляционных матриц, учитывающей конечность экспериментальных записей. Применение этого алгоритма к данным кальциевого имиджинга показало, что эффективная размерность обратно связана со скоростью движения животного, снижаясь в периоды локомоции, что согласуется с усилением синхронной активности нейронов гиппокампа при движении. Исследование внутренней размерности методом ближайших соседей подтвердило, что активность гиппокампальных популяций лежит на сильно искривлённом низкоразмерном многообразии, размерность которого остаётся почти постоянной во времени, в отличие от линейных мер. Обнаружено степенное масштабирование внутренней размерности с числом регистрируемых нейронов, которое не объясняется структурой временного ряда и является характеристикой коллективной нейронной активности.

Для учёта темпоральной структуры активности, не используемой методами снижения размерности, разработан мультиплексный граф рекуррентности популяции~--- авторский способ агрегации индивидуальных рекуррентных графов нейронов в единое представление; насколько известно, метод впервые применён к сырым сигналам кальциевого имиджинга \textit{in~vivo}. Из одного и того же графа восстанавливаются как геометрия среды (медианная ошибка декодирования положения животного около 9~см при стороне арены 44~см), так и динамические режимы коллективной активности, разделяемые по степени регулярности и связанные со скоростью движения. Тем самым рекуррентный подход дополняет анализ эффективной размерности, давая геометрию состояний поверх динамической характеристики.

На уровне индивидуальных нейронов разработан информационно-теоретический фреймворк INTENSE для выявления и количественной оценки нейронной селективности в данных кальциевого имиджинга при свободном поведении. Метод работает непосредственно с кальциевым сигналом без деконволюции, обрабатывает переменные различных типов и включает ускоренную процедуру статистического тестирования на основе циклических сдвигов с контролем автокорреляции. При применении к гиппокампальным записям выявлены разнообразные типы селективных нейронов; показано, что примерно треть наблюдаемой мультиселективности объясняется поведенческой избыточностью, а замена взаимной информации на линейную корреляцию приводит к потере более двух третей значимых ассоциаций. Применение INTENSE к искусственной РНС подтвердило функциональные роли динамически определённых популяций и выявило тонкую внутреннюю структуру, недоступную при популяционном анализе.

Дополняющий INTENSE подход реализован во фреймворке MANGO, который впервые обеспечил безградиентную максимизацию активации нейронов в импульсных нейронных сетях на основе оптимизации с разложением тензорного поезда. Если INTENSE характеризует селективность со стороны статистической зависимости между активностью и внешними переменными, то MANGO~--- со стороны генерации оптимального стимула. Анализ выявил нетривиальную двухфазную динамику формирования селективности в импульсных сетях, отличающуюся от монотонного процесса в классических ИНС, и установил параллели с критическими периодами в биологических нейронных сетях. Последующее применение к выходам импульсной сети инструментов снижения размерности и информационного анализа (DRIADA, INTENSE) выявило двухфазную траекторию в информационной плоскости глубоких слоёв и усиление связи индивидуальной селективности нейронов с глубиной сети. Это демонстрирует субстрат-независимость разработанного инструментария: методы, созданные для нейрофизиологических данных, содержательно применимы и к искусственным нейронным сетям.

Все разработанные методы интегрированы в открытый программный пакет DRIADA, объединяющий инструменты анализа нейронной активности на индивидуальном, популяционном и сетевом уровнях в рамках единой модели данных. Пакет использовался для получения результатов, представленных во всех разделах данной диссертации, обеспечивая согласованность форматов данных и статистических процедур. Сопутствующий инструмент автоматического контроля качества выделенных компонент кальциевого имиджинга обеспечил воспроизводимую предобработку записей для всех нейрофизиологических разделов работы.

Объединение индивидуального и популяционного уровней анализа единым информационно-теоретическим критерием показало, что популяционный код гиппокампа сильно распределён: в формирование коллективных переменных вовлечена почти вся популяция, тогда как вклад отдельного нейрона умеренный. При этом поведенческая селективность нейрона значимо положительно связана с его вовлечённостью: поведенчески селективные нейроны вовлечены в код сильнее, хотя он формируется широкой смесью селективных и неселективных клеток.

Перспективными направлениями дальнейших исследований являются: разработка подходов, устойчивых к гетерогенности коллективных переменных, когда отдельные моды (как пространственное кодирование в гиппокампе) доминируют над остальными; а также применение фреймворка MANGO к биологическим нейронам \textit{in~vivo}.
